\section{Hardware Software Implementation}
\label{sec:hardware_software_implementation}

\subsection{Hardware Implementation}

The developed prototype consists of a compact autonomous mobile robot capable of line following and obstacle avoidance. The system integrates sensing, control, actuation, and power management into a single embedded platform. Figure~\ref{fig:robot} shows the completed robot prototype used during implementation and testing.

The Arduino Uno microcontroller serves as the central controller of the robot. It continuously receives data from two infrared (IR) sensors and two HC-SR04 ultrasonic sensors, processes the sensor information, and generates motor control commands. The IR sensors are mounted at the front of the robot to detect the black line, while the ultrasonic sensors continuously monitor the distance to obstacles during navigation.

An L293D motor driver is used to control two DC geared motors. The motor driver receives PWM signals from the Arduino and regulates the speed and direction of both motors. To compensate for mechanical differences between the motors, calibrated PWM values were applied to obtain stable straight-line movement.

The complete system is powered by a 12\,V battery. The motor driver supplies a regulated 5\,V output to power the Arduino Uno and the connected sensors. This configuration provides a simple and reliable power distribution while maintaining stable operation of both the electronic and mechanical components.

The electronic components are mounted on a lightweight wooden chassis using a breadboard for rapid prototyping. This modular arrangement simplifies wiring, maintenance, and future hardware modifications.

\begin{table}[ht]
\caption{Hardware Components}
\centering
\begin{tabular}{lll}
\hline
\textbf{Component} & \textbf{Qty} & \textbf{Purpose} \\
\hline
Arduino Uno & 1 & Main controller \\
L293D Motor Driver & 1 & Motor control \\
DC Motors & 2 & Robot movement \\
IR Sensors & 2 & Line detection \\
HC-SR04 Sensors & 2 & Obstacle detection \\
Breadboard & 1 & Circuit prototyping \\
12\,V Battery & 1 & Power supply \\
Wooden Chassis & 1 & Mechanical support \\
\hline
\end{tabular}
\label{tab:hardware}
\end{table}

\begin{figure}[ht]
\centering
\includegraphics[width=0.48\linewidth]{robot.png}
\caption{Final hardware prototype of the autonomous line-following and obstacle-avoidance robot.}
\label{fig:robot}
\end{figure}

\subsection{Software Implementation}

The control software was developed using the Arduino Integrated Development Environment (Arduino IDE) in C++. A modular programming approach was adopted to separate sensor acquisition, motor control, obstacle detection, and navigation into independent software functions. This improves readability, simplifies debugging, and allows future system extensions.

During initialization, the Arduino configures all digital input and output pins, initializes the infrared and ultrasonic sensors, and enables the motor driver. After initialization, the software repeatedly executes the main control loop throughout the robot's operation.

The first stage of the control algorithm continuously checks the ultrasonic sensors for nearby obstacles. When an obstacle is detected within the predefined safety distance, the robot immediately stops and confirms the measurement before executing the obstacle avoidance routine.

If no obstacle is detected, the robot performs line following using the two infrared sensors. Based on the sensor readings, the controller determines whether the robot should move forward, turn left, or turn right. Independent PWM control of the two motors enables smooth steering corrections while following the track.

To improve navigation reliability, the software stores the previous movement direction. If both infrared sensors temporarily lose the line, the robot searches in the last known direction until the track is detected again. After avoiding an obstacle, the robot continues searching until the line is reacquired before resuming normal navigation.

Overall, the implemented software provides reliable autonomous operation while maintaining a computationally efficient control strategy suitable for embedded systems.

\subsection{Hardware--Software Integration}

The hardware and software operate together as a closed-loop embedded system. Sensor information collected by the infrared and ultrasonic sensors is processed by the Arduino Uno, which determines the appropriate movement commands. These commands are transmitted to the L293D motor driver, which controls the speed and direction of the two DC motors.

The continuous interaction between sensing, processing, and actuation enables the robot to follow a predefined path, detect obstacles, execute avoidance manoeuvres, and return to the original track without external intervention. The implementation demonstrates that reliable autonomous behaviour can be achieved using low-cost hardware combined with an efficient embedded control algorithm.
